% Generated by tex/build_swebok_structure.py from tex/chapters/ch01/00_chapter_opening.tex
\section{第01章 ソフトウェア要求}
\noindent\textbf{この資料の主張}

ソフトウェア要求は、単なる「欲しい機能の一覧」ではない。要求は、現実世界の問題を解くためにソフトウェアが満たすべき条件、能力、制約を言葉やモデルで表したものである。要求が曖昧なまま設計・実装・テストへ進むと、後で直す範囲が大きくなる。したがって、要求を「引き出す」「分析する」「仕様として書く」「妥当か確認する」「変化を管理する」という一連の活動として扱うことが重要である。

\MiniBox{この資料の読み方}{専門用語は原則として日本語で示し、必要に応じて英語を括弧内に併記する。英語名は、元資料やツール上の名称を確認したいときの手がかりとして残す。初学者が前提知識なしに読めるように、各節では「何のための概念か」「実務では何を見ればよいか」を重視する。}
